%%%%%%%%%%%%%%%%%%%%%%%%%%%%%%%%%%%%%%%%%
% Journal Article
% LaTeX Template
% Version 2.0 (February 7, 2023)
%
% This template originates from:
% https://www.LaTeXTemplates.com
%
% Author:
% Vel (vel@latextemplates.com)
%
% License:
% CC BY-NC-SA 4.0 (https://creativecommons.org/licenses/by-nc-sa/4.0/)
%
% NOTE: The bibliography needs to be compiled using the biber engine.
%
%%%%%%%%%%%%%%%%%%%%%%%%%%%%%%%%%%%%%%%%%

%----------------------------------------------------------------------------------------
%	PACKAGES AND OTHER DOCUMENT CONFIGURATIONS
%----------------------------------------------------------------------------------------

\documentclass[
	a4paper, % Paper size, use either a4paper or letterpaper
	10pt, % Default font size, can also use 11pt or 12pt, although this is not recommended
	unnumberedsections, % Comment to enable section numbering
	twoside, % Two side traditional mode where headers and footers change between odd and even pages, comment this option to make them fixed
]{LTJournalArticle}

\addbibresource{sample.bib} % BibLaTeX bibliography file

\runninghead{Processing Personal Data under the EU AI Act} % A shortened article title to appear in the running head, leave this command empty for no running head

\footertext{AI: Regulation \& Law} % Text to appear in the footer, leave this command empty for no footer text

\setcounter{page}{1} % The page number of the first page, set this to a higher number if the article is to be part of an issue or larger work

%----------------------------------------------------------------------------------------
%	TITLE SECTION
%----------------------------------------------------------------------------------------

\title{Processing Personal Data under the EU AI Act:\\Interaction and Conflict with the GDPR} % Article title, use manual lines breaks (\\) to beautify the layout

% Authors are listed in a comma-separated list with superscript numbers indicating affiliations
% \thanks{} is used for any text that should be placed in a footnote on the first page, such as the corresponding author's email, journal acceptance dates, a copyright/license notice, keywords, etc
\author{%
	Marlon Müller\textsuperscript{1}
}

% Affiliations are output in the \date{} command
\date{\footnotesize\textsuperscript{\textbf{1}}Technical University of Munich, Germany}

%----------------------------------------------------------------------------------------

\begin{document}

\maketitle % Output the title section

%----------------------------------------------------------------------------------------
%	ARTICLE CONTENTS
%----------------------------------------------------------------------------------------

\section{Introduction}

The training of Artifical Intelligence (AI) systems relies heavily on the processing of large amounts of data. In practice a large share of this data is personal data. This overlap between AI and personal data processing brings two major 
regulatory instruments of the European Union (EU) into contact: the General Data Protection Regulation (GDPR) and the Artificial Intelligence Act (AIA).\\
Both regulatory frameworks have the same goal of protecting fundamental rights and freedoms of individuals, and they both apply simultaneously to to providers of high-risk AI systems that process personal data.
This relationship is conceptualized as an additive one in Art. 2(7) AIA~\cite{euaiact2024}. They also follow a different regulatory logic, with the GDPR protection the fundamental right of privacy and the AIA acting as a product safety regulation for AI systems.\\
This paper analyzes how this co-regulation affects providers and argues that the relationship is more complex than the AIA's additive model suggests. For this purpose, the paper first establishes a role mapping and 
will then examine the parallel obligations that arise from the two frameworks. Finally, it will show that the debiasing exception in Art. 10(5) AIA requires an interaction between the two frameworks.\\

\section{Two Frameworks, One Actor}

In both regulations there are distinct roles that entities can take on and that have different legal obligations. The GDPR distinguishes between data controllers and data processors, while the AIA distinguishes between providers and deployers of AI systems.
In practice the roles of the two frameworks do not map onto each other statically, even though they often overlap if personal data that falls under the GDPR is involved.\\
During the development and training of an AI system the provider can usually be considered to be the data controller, as it determines the purposes and means of the processing of personal data.
Even though the provider can still give instructions to the deployer during the deployment phase on how to use the system and how to handle personal data, these constraints are merely technical and are not sufficient to make the provider a data processor during
deployment. This role will be taken on by the deployer instead, as he keeps control over the purposes of the processing of personal data and, most crucially, can control who is a data subject, usually customers of the deployer, of that processing.
If the inference does not run on the deployer's infrastucture and personal data flows to the server the AI system is hosted on, the natural or legal person that operates that server, which is either the provider, importer or distributor of the AI system,
can be considered a data processor during inference. If the system is self-hosted by the deployer, the provider exits the processing chain entirely and is not subject to GDPR after deployment.~\cite{euaiact2024, eugdpr2016}\\
While the mapping of these roles is grounded in the legal definitions of the two frameworks and the practical realities of AI development and deployment, there may nevertheless be edge cases or legal and contractual constellations in which this mapping does not hold.\\
\\
One important gap of the overlap of these roles can be found whenever the provider of a continuous learning system feeds inference data back into retraining.
In this case, the provider may transition from processor to controller for that secondary processing, as it determines a new purpose for its own benefit. 
This practice could conflict with GDPR Articles 5(1)(b) and 6~\cite{eugdpr2016}, but the AIA addresses it only obliquely through Article 15 on feedback loops, without adressing or resolving the GDPR consequences~\cite{euaiact2024}.\\

\section{Parallel Obligations}
In Article 2(7) the AIA states that the obligations of the GDPR apply in addition to the obligations of the AIA. This means that providers and deployers of high-risk AI systems that process personal data are subject to both sets of obligations simultaneously.
This creates a complex regulatory landscape for these entities, as they must navigate and comply with the requirements of both frameworks. The GDPR has very specific obligations regarding data protection, such as data minimization, purpose limitation, and the rights of data subjects, while the AIA has its own set of obligations related to the safety and transparency of AI systems.\\
In the following, the paper will analyze some of the most important obligations that arise from this co-regulation and wether they are complementary or conflicting.

\subsection{Transparency Obligations}
Both the GDPR and the AIA have certain transparency obligations, but they differ significantly in their purpose and scope. The GDPR defines its transparency obligations for data controllers in Articles 13 and 14.
Main purpose of these obligations is to inform the data subjects about the processing of their personal data, including the purposes of the processing, the categories of data being processed as well as the recipients and time of storage of the data.
It is notable that the GDPR does not require data controllers to disclose any information about the inner workings of the processing itself, such as the algorithms used or the logic behind the processing.~\cite{eugdpr2016}\\
The AIA, on the other hand, has transparency oblications that are focused on the interpretability and explainability of the AI system itself. These obligation are outlined in Article 13(1) and Article 13(3)(f) and requires providers to provide instructions for use with their system. 
However, these instructions for use are to be provided to the deployer and not to the data subjects, and they are not required to include any information about the processing of personal data specifically. 
The only transparency requirement, directly targeting natural persons, and therefore potentially data subjects, is the disclosure of the fact that they are interacting with an AI system, as defined in Article 50 AIA.~\cite{euaiact2024}\\
While both frameworks have transparency obligations, they are not necessarily complementary, as they target different audiences and have different purposes. As they also don't cause any direct conflicts, these obligations can be seen as parallel and independent of each other.

\subsection{Data Governance}
The governance of data is a crucial aspect of the AIA, especially to ensure the quality and representativeness of the data used for training and testing AI system. Apart from data quality Article 10(2)(b) AIA also demands that providers consider the origin of the data
as well as the original purpose of the collection of personal data~\cite{euaiact2024}. Even though this clause references the original purpose of the collection, it does not explicitly mandates the compliance with the purpose limitation principle 
established in Article 5(1)(b) GDPR. Not only does the AIA fail to explicitly require compliance with the GDPR's purpose limitation principle, it requires providers to ensure a representative and bias free dataset, which, in practice,
requires large datasets~\cite{Samuel2021Distributional} as well as the collection of additional data that may not be necessary for the original purpose of the collection, and therefore conflicts with the GDPR's principle of data minimization in Article 5(1)(c)~\cite{eugdpr2016}.

\subsection{Right to be Forgotten}
The right to be forgotten is established in Article 17 GDPR and allows data subjects to request the deletion of their personal data under certain conditions. This right is regularly confirmed and strengthend by the Court of Justice of the European Union (CJEU) and other
national courts~\cite{googlespain2014,Kazemi2018}. Personal data that is used for training an AI system becomes an integral part of the AI model and is not stored in a traditional sense, but rather encoded in the parameters of the model and the influence of a single data
point on the model is not easily identifiable. In cases where it is possible to extract the training data from the model in a way that could allow the identification of a data subject the European Data Protection Board (EDPB) has clarified that the right to be 
forgotten applies to the trained model as well~\cite{edpb}. In practice this requires the AI model to unlearn this data point, which is technically possible but may not always be feasible~\cite{Velli2025,JuliussenRuiJohansen2023}.
Even though the right to be forgotten is a fundamental right under the GDPR and compliance with it is a legal obligation for data controllers, the AIA does not explicitly ensures that providers of AI systems have the technical means to unlearn data points from their
model or any other way to comply with the right to be forgotten. This leaves a gap between a significant obligation under the GDPR and the requirements of the AIA.\\
\\
In general these obligations show that the co-regulation of the GDPR and the AIA are additive in some cases, but they also have significant regulatory gaps and structural tensions between the compliance with both frameworks.\\

\section{The Debiasing Exception}

So far it could be established that, while there are some structural gaps and tensions between the two frameworks, there is no substantial interaction between the GDPR and the AIA. However, this co-regulation paradigm changes in
Art. 10(5) AIA, which establishes an exception of the GDPR for the purpose of debiasing AI systems. In general Art. 10(2)(f) AIA demands that providers of high-risk AI systems ensure that the data used for training is free of any biases
that could lead to discriminatory outcomes~\cite{euaiact2024}, especially these biases could be related to disablity, religion, racial origin or sexual orientation.

\subsection{Proxy characteristics}

Even data that seems unrelated to any of the aforementioned protected characteristics could classify as a proxy for those characteristics. To give a concrete example, the Police Crime Statistics in Germany (PCS) shows that crime suspects
with a non-German citizenship are overrepresented when compared to their share of the general population~\cite{pks2025}. A closer look into the data reveals that this overrepresentation vanishes when the variables are controlled for factors 
such as age, gender, place of living and socio-economic status~\cite{pksIfo}. This means that the non-German citizenship is not a predictor for criminal behaviour, in fact those two variables are not even correlated\footnote{There are even more reasons why the PCS is not suitable to establish a causal relationship between having a non-German citizenship and criminal behaviour \cite{pksIfo,Integration}.}, 
but the non-German citizenship acts as a proxy for the other factors, which are the actual predictors for criminal behaviour, and vice versa.

\subsection{No debiasing without personal data}

The example of the PCS shows that a provider of an AI system that uses crime data for training needs to ensure that these biases are removed from the training data. Without access to the migration status of crime suspects,
it would be impossible to identify and remove the bias from the data~\cite{VANBEKKUM2025106115}. Unfortunately, the processing of special categories of personal data, including the aforementioned disability, religion, racial origin and sexual orientation,
is generally prohibited under Art. 9(1) GDPR, unless one of the exceptions in Art. 9(2) applies. Art. 9(2)(g) GDPR establishes the exception for processing these special categories of personal data where it is necessary for reasons of substantial
public interest and is based on Union law.~\cite{eugdpr2016} The AIA is a Union law under this Article~\cite{VANBEKKUM2025106115} and allows the processing of special categories of special categories of personal data for
the pupose of debiasing AI systems in Art. 10(5)~\cite{euaiact2024}. The public interest of this exception, which is a requirement of Art. 9(2)(g) GDPR, was justified in Recital 70 of the AIA.\\
\\
The debiasing exception in Art. 10(5) AIA shows that the relationship between the GDPR and the AIA is not purely additive, but that it constitutes a significant interaction between the two frameworks.
The exception is also necessary to ensure that high-risk AI systems can be trained on bias-free data, which is crucial to ensure non-discriminatory outcomes and to protect the fundamental rights of individuals.\\

\section{Conclusion}

The relationship of the GDPR and the AIA operates in three different modes. First, there are additive obligations, like the transparency obligations, that apply in parallel. On the other hand, 
the priniciple of data minimization of the GDPR and the data governance obligations of the AIA create competing requirements. Additionally, the practical consequences of the right to be forgotten in AI systems show a significant gap between the GDPR and the AIA. 
Finally, Art. 10(5) AIA establishes a de-biasing exception that allows the processing of special categories of personal data, which is normally prohibited under Art. 9(1) GDPR.
Due to the three different modes of co-regulation providers of high-risk AI systems cannot simply apply a single compliance logic, but need to determine the mode of interaction provision by provision.\\

%----------------------------------------------------------------------------------------
%	 REFERENCES
%----------------------------------------------------------------------------------------

\printbibliography % Output the bibliography

%----------------------------------------------------------------------------------------

\end{document}
